\section{The two-way pilot}
\label{app:pairs}

Two copies that read each other can end in one of three ways: each
keeps its own assignment, the two exchange assignments, or both settle
on one rule. Which of these happens cannot be read from either member
alone, because a pair can converge on either member's rule, so we score
the two members' answers as a pair. The pilot tested three two-way
conditions, defined below: mutual reading of fixed memories, the live
loop, and fixed memory plus a weak live read. Mutual reading of fixed
memories exchanged the assignments. Only the live loop pulled pairs
onto one member's rule more often than independent choices predict,
most clearly on the current rule and only while the link was open
(Section~\ref{sec:twoway}). Fixed memory plus a weak live read behaved
much like mutual reading of fixed memories at $w=2$.

\paragraph{Design.}
In \emph{mutual reading of fixed memories}, each member reads the
other's fixed memory (its cached card, note, and reflection prompt) at
weight $w$. Neither memory changes, so neither member can respond to
the other, and the condition amounts to two one-way links running side
by side. In the \emph{live loop}, each reads the other's whole cache at
weight $w$, including the reflection tokens as they are written, one
decoding step behind, so each can respond to what the other writes. In
\emph{fixed memory plus a weak live read}, each member weights the
partner's fixed memory by $w_P=2$ and the tokens the partner writes
afterwards, its continuation, by $w_C=.1$. This condition asks whether a
small channel for responding, laid over strong fixed-memory reading,
changes what strong fixed-memory reading does on its own. The live loop
and fixed memory plus a weak live read were also run one-way, with only
A reading B; we call the first of these the \emph{one-way live loop}.
After 48 reflection tokens, the pair's state is saved and copied into two
independent branches: A is questioned in one while B keeps reflecting,
and B in the other. A pair outcome combines the two answers within each
option ordering, so agreement means that the members gave the same
answer when asked separately, not at the same moment.

\paragraph{Weights.}
The live loop and fixed memory plus a weak live read both ran at the
weakest setting in their grids, because every stronger setting cost too
much capability. The weights came from a separate 40-episode two-way
strength pilot that looked only at capability and answer format. A
setting qualified if, in the two-way condition, both members lost at
most five points of capability and met the format criteria, and we
selected the strongest qualifying setting of each condition. At the
selected settings, capability fell by 2.50 points for A and 3.75 for B
in the live loop at $w=.5$, and by 3.75 and 3.12 points in fixed memory
plus a weak live read at $w_C=.1$. Every stronger setting took at least
one member beyond the tolerance: the larger of the two declines was 6.25
and 8.12 points for the live loop at $w=.55$ and $.6$, and 5.62 and
12.50 points for fixed memory plus a weak live read at $w_C=.3$ and
$.5$. Whether a stronger live loop would do more than couple current
rules is therefore untested.

\paragraph{Pair outcomes.}
On a given question, a pair \emph{keeps} (each member selects its own
rule), \emph{A takes B's} (A selects B's rule while B keeps
it), \emph{B takes A's} (the converse), \emph{swaps} (each selects the
other's rule), agrees on a \emph{third} rule (both select
Robin's or the unused one), or falls into \emph{other} (any remaining
combination). In the two taking outcomes both members hold one member's
rule; their sum is the one-side-wins rate $W$. Each episode
contributes the mean over its two orderings, and the episode is the unit
for every interval.

\paragraph{Excess convergence.}
Two members that are each pulled toward the partner, independently of
each other, still produce one-side-wins pairs by chance. The excess
removes that expectation. With $p_{A\to B}$ and $p_{B\to A}$ the members'
marginal rates of choosing the partner's rule, and $p_{A\to A}$ and
$p_{B\to B}$ their rates of keeping their own, each averaged with equal
episode weight within a condition,
\begin{equation}
E=W-\left(p_{A\to B}\,p_{B\to B}+p_{A\to A}\,p_{B\to A}\right).
\label{eq:pair-excess}
\end{equation}
Pooling episodes can create excess of its own, for example when one
member's rule is more salient in some episodes than in others. The
primary comparison therefore subtracts the excess of a control in which
the members read the same fixed content but cannot respond: mutual
reading of fixed memories at the matching weight. Each bootstrap draw
resamples episodes and recomputes both conditions' marginals. The
pilot's protocol set the condition for moving to a confirmatory study:
a difference of at least 10 percentage points on the assigned rule, with
a 90\% interval excluding zero. The same difference on the current
rule, and on the assigned rule after the link was cut, were
prespecified secondary comparisons. Table~\ref{tab:pair-gate} gives all
six comparisons. Neither condition reached the 10-point threshold, so
the two-way study did not go on to a confirmatory sample.

\begin{table}[htbp]
\centering\small
\caption{Prespecified comparisons in the two-way pilot ($N=80$):
excess convergence (Eq.~\ref{eq:pair-excess}) in the condition, in
its control (mutual reading of fixed memories, at $w=.5$ for the live
loop and at $w=2$ for fixed memory plus a weak live read), and their
difference, in percentage points with 90\% episode-bootstrap intervals.
\emph{Fixed $w=2$ $+$ weak live}: fixed memory plus a weak live read.
The first row of each condition is its primary comparison, which
required a difference of at least 10 points; neither condition reached
it. With the link cut, neither condition left any excess.}
\label{tab:pair-gate}
\begin{tabularx}{\linewidth}{@{}Ylrrr@{}}
\toprule
Condition & Question, link & Excess & Control & Difference [90\% CI] \\
\midrule
Live loop, $w=.5$ & Assigned rule, open (primary) & 6.875 & .141 & $6.734\ [2.617,10.906]$ \\
 & Current rule, open & 12.430 & .312 & $12.117\ [7.391,16.469]$ \\
 & Assigned rule, cut & .000 & .000 & $.000\ [.000,.000]$ \\
\addlinespace
Fixed $w=2$ $+$ weak live & Assigned rule, open (primary) & .328 & .156 & $.172\ [.000,.445]$ \\
 & Current rule, open & 1.625 & .859 & $.766\ [.234,1.430]$ \\
 & Assigned rule, cut & .000 & .000 & $.000\ [.000,.000]$ \\
\bottomrule
\end{tabularx}
\end{table}

\paragraph{What the complete outcomes show.}
Table~\ref{tab:pair-outcomes} lists every condition on both questions.
Three patterns bear on the main text. First, fixed and live reading move
pairs in different directions. Mutual reading of fixed memories
exchanges assignments at $w=2$ (94.38\% swaps on the assigned rule)
and leaves most pairs unchanged at $w=.5$. At that same weight the live
loop moves pairs onto one member's rule instead, and swaps stay
rare. Second, no pair in any condition agreed on a third rule, so
convergence, where it occurred, was always onto one member's
assignment. Third, cutting the link returned every live-loop pair to its
own assignments on both questions. Mutual reading of fixed memories at
$w=2$ lost its swaps but kept a residue on the current rule: in
41.2\% of pairs one member still named the other's rule, the
pairwise counterpart of the one-way residue in
Section~\ref{sec:moments}.

% Complete pair outcomes of the two-way pilot; episode means, copied without recalculation.
\begin{table}[htbp]
\centering\small
\setlength{\tabcolsep}{3pt}
\caption{Complete pair outcomes in the two-way pilot, in percent of pairs
($N=80$ episodes per row; each episode contributes the mean of its two
orderings), on the questions about the assigned and the current
rule. \emph{Keep}: each member selects its own rule; \emph{A
takes B's}: A selects B's rule while B keeps it; \emph{B takes
A's}: the converse; \emph{Swap}: each selects the other's rule;
\emph{Other}: any remaining combination. No pair agreed on a third
rule in any condition. A's marginal rate of choosing B's rule is
the sum of the A-takes and Swap columns. \emph{Mutual fixed}: mutual
reading of fixed memories. \emph{Fixed $w=2$ $+$ weak live}: fixed memory
plus a weak live read, which weights the partner's fixed memory by
$w_P=2$ and the tokens the partner writes afterwards by $w_C=.1$.
\emph{One-way live loop}: the live loop with only A reading B.}
\label{tab:pair-outcomes}
\begin{tabularx}{\linewidth}{@{}Yrrrrrrrrrr@{}}
\toprule
& \multicolumn{5}{c}{Assigned rule} & \multicolumn{5}{c}{Current rule} \\
\cmidrule(lr){2-6}\cmidrule(l){7-11}
Condition & Keep & \shortstack{A takes\\B's} & \shortstack{B takes\\A's} & Swap & Other & Keep & \shortstack{A takes\\B's} & \shortstack{B takes\\A's} & Swap & Other \\
\midrule
No link & 100.00 & 0.00 & 0.00 & 0.00 & 0.00 & 100.00 & 0.00 & 0.00 & 0.00 & 0.00 \\
\addlinespace
\multicolumn{11}{@{}l}{\emph{One-way: only A reads B}} \\
One-way live loop, $w=.5$ & 91.88 & 8.12 & 0.00 & 0.00 & 0.00 & 85.00 & 15.00 & 0.00 & 0.00 & 0.00 \\
Fixed $w=2$ $+$ weak live & 3.75 & 96.25 & 0.00 & 0.00 & 0.00 & 6.88 & 93.12 & 0.00 & 0.00 & 0.00 \\
\addlinespace
\multicolumn{11}{@{}l}{\emph{Two-way, link open}} \\
Mutual fixed, $w=.5$ & 93.12 & 1.25 & 5.62 & 0.00 & 0.00 & 91.25 & 2.50 & 6.25 & 0.00 & 0.00 \\
Mutual fixed, $w=2$ & 0.00 & 2.50 & 3.12 & 94.38 & 0.00 & 0.00 & 6.88 & 6.25 & 86.88 & 0.00 \\
Live loop, $w=.5$ & 53.75 & 17.50 & 25.62 & 1.88 & 1.25 & 43.12 & 26.25 & 27.50 & 2.50 & 0.62 \\
Fixed $w=2$ $+$ weak live & 0.00 & 4.38 & 3.75 & 91.88 & 0.00 & 0.00 & 10.00 & 8.12 & 81.88 & 0.00 \\
\addlinespace
\multicolumn{11}{@{}l}{\emph{Two-way, link cut before the question}} \\
Mutual fixed, $w=.5$ & 100.00 & 0.00 & 0.00 & 0.00 & 0.00 & 98.75 & 0.00 & 1.25 & 0.00 & 0.00 \\
Mutual fixed, $w=2$ & 96.88 & 0.00 & 3.12 & 0.00 & 0.00 & 58.75 & 18.75 & 22.50 & 0.00 & 0.00 \\
Live loop, $w=.5$ & 100.00 & 0.00 & 0.00 & 0.00 & 0.00 & 100.00 & 0.00 & 0.00 & 0.00 & 0.00 \\
Fixed $w=2$ $+$ weak live & 96.88 & 0.00 & 3.12 & 0.00 & 0.00 & 62.50 & 18.12 & 19.38 & 0.00 & 0.00 \\
\bottomrule
\end{tabularx}
\end{table}


\paragraph{Does a responsive partner pull harder?}
In the one-way live loop, B does not read A, so B's reflection runs on
regardless of what A writes. In the live loop B also reads A and can
respond. A post hoc comparison of the two at $w=.5$ asks whether a
partner that can respond pulls the receiver further. In this pilot it
did. Letting B also read A raised the share of A's answers naming B's
rule from 8.1\% to 19.4\% on the assigned rule (11.2 points,
90\% CI $[5.0,18.1]$) and from 15.0\% to 28.8\% on the current rule
(13.8 points, 90\% CI $[8.1,20.0]$). Each rate is A's marginal rate of
choosing B's rule, the sum of the A-takes and swap columns of
Table~\ref{tab:pair-outcomes}, and the intervals are episode-bootstrap
intervals for the paired difference.

A prespecified test of the same comparison went the other way. It ran
as an 80-episode signal pilot under the second protocol and compared
the live loop at $w=.7$ with the one-way live loop at the same weight,
on A's source score for the current rule, $M_{\mathrm{NOW}}$, alone. The
weight $.7$ was the strongest weight for reading the partner's evolving
cache that had passed the capability limit in that protocol's strength
pilot, in which only A read B. The
difference was $-4.79$ nats, 90\% CI $[-8.31,-1.18]$, opposite to the
predicted direction. Two features of that test limit what it can show.
In the signal pilot itself, capability fell by 9.4 points one-way and
14.4 points two-way, beyond the five-point tolerance, so both arms were
answering in a degraded state. And a score built from A alone counts
only movement toward B's rule, so a pair converging on A's rule would
lower it. For these reasons the two-way pilot set its weights so that
both members stayed within the capability tolerance, and it scored both
members, but whether either feature explains the reversal is untested. The two tests disagree, and
only a confirmatory sample could settle which holds.

Even the post hoc difference need not mean that B responded to what A
wrote while reflecting. With the link kept, A reads the question one
token per decoding step while B goes on reflecting
(Appendix~\ref{app:methods}). In the live loop B also reads A's cache,
which by then contains the question, so B's next tokens may reply to it,
for example by stating B's own rule, and A reads those tokens before
it answers. In the one-way live loop B cannot see the question. The
difference may therefore come from B's reaction to A's question at
answer time rather than from anything the pair built up during the
reflection. A run in which B cannot read A's question tokens would
separate the two.

The same comparison with fixed memory plus a weak live read found no
difference. There A already named B's rule in 96.2\% of answers
about its assigned rule when only A read B, and in 96.2\% when B
also read A; on the current rule the rates were 93.1\% and 91.9\%.

\paragraph{Fixed memory plus a weak live read.}
Adding a weak live read to strong fixed-memory reading did not turn
swaps into convergence. As $w_C$ goes to zero this condition reduces to
mutual reading of fixed memories at $w=2$, which is therefore its
control. With the link open, 91.9\% of pairs still swapped on the
assigned rule, against 94.4\% without the weak live read, and the weak
live read added less than one point of excess convergence on either
question (Table~\ref{tab:pair-gate}). After the cut the condition
behaved like its control as well. In 37.5\% of pairs one member still
named the other's current rule, against 41.2\% for mutual reading of
fixed memories alone, so the weak live read left no residue of its own.
